\subsection{Executable deployments in serQ v0.1.0}\label{L1:sec:language}

A serQ program declares memory pools, service stages, a workload and a
session or server block. Its executable definition is IR version 9.
The examples accompanying these notes use this release's syntax.
A FIFO stage serves work in arrival order, a PS stage shares capacity,
and a delay stage gives each session its own timer. For example:
\begin{verbatim}
stage prefill : fifo;
stage link : ps(1);
stage decode : ps(min(present, 16));
stage tool : delay;
\end{verbatim}
Work at these stages is measured in seconds at unit rate.
A step engine instead schedules tokens and charges time per iteration:
\begin{verbatim}
stage engine : step {
  budget B;
  serve admission;
  cost max(omega, tokens * a + b * attention);
  memory kv;
}
\end{verbatim}
Here \texttt{tokens} is the iteration's scheduled token count,
\texttt{attention} sums $n(K+n/2)$ over prefill chunks, and
\texttt{present} is the number at a PS station. The declaration can choose
admission order, decode first, exclusive prefill, or explicit ordering keys.
FIFO claims require the selected order, rather than following from chunking alone.

\paragraph{Client and server.} The workload describes arrivals and the turn/tool
loop; a client-side \texttt{request;} invokes a server block once.
The server specifies admission, prefill, transfer and decode.
\texttt{hidden o;} prevents scheduling keys from reading an output length
that the workload sampled but the scheduler does not know.
Expression and statement definitions use \texttt{def}, imported with
\texttt{use "file.sq";}. They expand to the executable IR.

\paragraph{Admission and lifetime.} A scoped \texttt{hold} reserves pool units,
runs its body, then frees the units, optionally retaining a cached prefix.
\texttt{reserve} states an admission gate larger than the initial allocation;
\texttt{reuse} limits how much cached state the turn can consume.
A pool can delegate admission to a step engine with \texttt{admit via engine;}.
Waiting prefixes then remain evictable until scheduled. A \texttt{lease}
keeps an allocation after the scope; \texttt{release} gives it back.
\texttt{transfer (X) from P to Q (T);} runs the link, records $T$ computed
tokens at the destination, and releases the source. Reserving the destination
before transfer can prolong the source lifetime; this is different from the
continuous tandem defined in \cref{L1:ex:program}.

\paragraph{Blocks and preemption.} Allocations round up to blocks and cached
prefixes round down. A preempted decode keeps generated tokens but may lose
their KV. At readmission, \texttt{computed} lets the program recover the
known prompt plus generated tokens, re-prefill the lost part, and decode only
the remaining output. Restarting from the original prompt changes the work.
Ending a session alone does not remove its reusable cached blocks.
An eviction key therefore needs an explicit end-of-program policy if it
values a finished session differently from a paused one.

\paragraph{What to execute.} The accompanying \texttt{programs/lecture\_pd.sq}
implements the continuous tandem of \cref{L1:ex:program}.
\texttt{programs/vllm.sq} and its imported library implement the release's
step-engine example. They are different models, not interchangeable
spellings. Run a program with \texttt{serq run file.sq --seed 1 --json};
\texttt{serq check file.sq} checks its syntax and executable structure.
Neither command measures a real serving system.
